.he
.bp76
.mt7
.mb9
.po8
.ll65
.ig
CHAPTER 19
S E T T I N G	 U P
.sp2
.pp5
The first step in setting up your Payroll System is to make
two copies of the distribution disks.
One distribution
disk contains the Startup programs, another contains the normal processing
programs, and a third contains the state and local tax table files and two
utility programs.
.he Operations: Chapter 19
.pp
To copy the disks
place a formatted disk containing the CP/M operating system,
and PIP
on drive A.  Place the (write protected) Startup distribution disk
on drive
B and hit the RESET switch on your machine.  Type the command,
.sp
.li
		A>PIP A:=B:*.*[ov]
.sp
and follow with RETURN.  When the PIP is complete, remove the distribution
disk from drive B, replacing it with a formatted disk that has the CP/M
system SYSGEN'ed onto it.
Enter CTRL-C or RESET the machine, and give
the command,
.sp
.li
		A>PIP B:=A:*.*[ov]
.sp
You now have a Startup Master Disk and a Startup Work Disk.  Repeat the
same process to create a Processing Master Disk and a Processing Work Disk.
Copy the Tax disk in the same way. Label one copy the Tax Master Disk and
the other copy the Payroll Data Disk.
The two work disks with the programs
will need CRUN2.COM PIP'ed onto them.  If you are operating
in single density you will probably not be able to fit CRUN2 and PIP on the
disks together.  Remove PIP from the work disks (if it is on them) and
PIP CRUN2 onto the work disks from another disk that does have PIP.  CRT.DEF
(if used) should also be PIP'ed onto these disks before you can begin.
.pp
If you can use the two disk configuration described in Chapter 7,
the programs on the Startup Disk and the Processing Disk may be combined
on one disk.
.pp
The second step in setting up the Payroll System is to define the
characteristics of your CRT to the Display Management System.
.pp
If you have the Hazeltine 1500 CRT, the Payroll system will run
immediately; you may erase all files with the filetype of DEF, and the
CRTDEF.INT program from your system master and work disks (but not the
Applewood Computers distribution disk!).  You can use the commands,
.sp
.li
		ERA *.DEF
	   and, ERA CRTDEF.INT
.pp
If you have one of the CRT devices listed below, you must use the CP/M
REN command to rename the DEF file corresponding to your CRT to the name
CRT.DEF:
.sp
.cp11
.li
       CRT			     DEF FILE
    --------------------------------------------------

    ADM 3A			     ADM3A.DEF
    BEEHIVE 150 		     B150.DEF
    BEEHIVE 152 		     B152.DEF
    BEEHIVE 157 		     B157.DEF
    CROMEMCO 3100		     CRO3100.DEF
    DYNABYTE NAKED TERMINAL	     DYNABYTE.DEF
    SOROC IQ 120		     SOROC120.DEF
    IMSAI VIOC			     VIOC.DEF
.sp
With the disk holding the DEF files on drive A, type the command:
.sp
.li
		A>REN A:CRT.DEF=A:nnnnnnnn.ttt
.sp
where nnnnnnnn.ttt stands for the filename and filetype (suffix)
of your CRT definition file.  Erase the remaining DEF files from
your disk one at a time (for example, ERA ADM3A.DEF, ERA B150.DEF, etc.).
Erase also the CRTDEF.INT program.
The Payroll system always expects to find the CRT.DEF file on
the currently logged disk drive.
.pp
If you have neither the Hazeltine 1500 device, nor one of
the CRT's listed above, see Chapter 32 for instructions on how to
create a CRT.DEF file that defines
your CRT to the Display Management System by running the
CRTDEF.INT program.
.pp
The third step in setting up the Payroll system is to erase the state
tax table files you do not need from your data disk and rename the local
tax table file.
The state
withholding tax table files (except for California) take the form,
.sp
.li
		PYSWTaa.101
.sp
where "aa" is the two character state abbreviation.  There are four
tax table files for California:
.sp
.cp4
.li
	PYCALS.101	     [S=Single]
	PYCALM.101	     [M=Married]
	PYCALH.101	     [H=Head of Household]
	PYCALX.100
.sp
Be sure to keep all four on the data disk if the California tables
are to be used.
.pp
The local tax table files have the filetype of LWT, and a filename that
indicates the locality (for example, NEWYORK.LWT).  Find the one for your
locality and rename it to PYLWT.101.  For example,
.sp
.li
		A>REN A:PYLWT.101=A:NEWYORK.LWT
.sp3
.cp12
.ce
MID-YEAR SETUP PROCEDURE
.pp
When you set up the Payroll System in mid-year, after the other setup
steps are complete, you must run the History Update program.  This
program can be run by choosing it from the menu.  You can run it
as often as necessary until the first payroll application has taken
place.	After the first application, the History Update program will
not run.
.pp
When you select the History Update program from the menu the Update
Employee History entry screen appears with the cursor waiting at
the Employee Number field.  You can enter the number of the employee
whose history you wish to update (all employees must be updated before
the first application), or give one of the following commands:
.sp
.li
	  CTRL-N: Get the NEXT record on the file;  or
		  the first if none have been accessed.
	  CTRL-B: Go BACK to the previous record; or to
		  the last if none have been accessed.
	  ESCAPE: End the Employee History Update section
		  of the program and go on to the remaining
		  section.
.sp
After accessing an employee record, move the cursor to the first
request by hitting RETURN.  Type in the information and then
hit ESCAPE to return the cursor to the Employee Number field
where you can call up the record for the next employee.  ESCAPE
again (form the Employee Number field) ends this section of the
program, begins a short processing program, and then
brings up the "mini" menu shown in Figure 19.1.
The standard commands, CTRL-B, CTRL-S, AND CTRL-R (go back a field,
save data to disk, and refresh a garbled DMS screen) can be given
at any point (although CTRL-B means something different when issued
from the Employee Number field).
.sp
.cp10
.li
	      Y O U R	 C O M P A N Y	  N A M E
			PAYROLL SYSTEM
  00/00/00	     HISTORY UPDATE MENU


	      1. UPDATE EMPLOYEE HISTORY TOTALS
	      2. UPDATE COMPANY LIABILITIES, VACATION, ETC.
	      3. UPDATE PAYMENTS HISTORY

		  ENTER SELECTION NUMBER: [ ]
.pp
The Update Employee History Totals entry screen is an exact reproduction
of the Update Employee History screen, except that there is no Employee
Number field since this screen presents the totals of the individual
employees.  You can make changes if necessary, but if the sum of the
individual totals does not equal the actual company totals you should
make sure that the individual employee information has been entered
correctly.  Press ESCAPE to end this section of the program and return
to the History Update Menu.
.pp
Menu choices 2 and 3 are explained in Chapter 7.  To enter the
information requested by those menu items, simply choose them
from the menu and type in the data.  Hit ESCAPE to return to the
History Update Menu.  ESCAPE from the menu brings up the Payroll
System menu.
.pp
If you want to change historical information for individual employees
or enter information for the first time for a newly added employee
(before the first payroll application), you would run the program
again.	Of course, if you do so you may have to adjust the new totals
through the Update Employee History Totals program again.
.he
.bp
.ig
CHAPTER 20
S E T T I N G	 S Y S T E M	P A R A M E T E R S
.sp2
.pp5
At first time startup, after you type CRUN2 PY to call up the system
the system menu appears.
The program asks you if you want to use the default Payroll System or
answer the full range of Parameter Entry requests (QUICK PARAMETER ENTRY
or PARAMETER ENTRY, respectively).
.he Operations: Chapter 20
.pp
If you choose the default Payroll System (menu item 10)
a limited parameter entry screen
appears, allowing you to enter your company name, address, and government
ID numbers. You also enter the pay interval you want to use (the default
system allows only one pay schedule), and the last day of the previous
pay period.  When you finish answering these requests and hit ESCAPE, the
startup sequence continues.  It creates the default account file,
the default System Deductions and Earnings file, the
Sort Parameter file, and then presents the system menu.  Once your
system is set up, to examine or change the parameter values
choose the Parameter Entry program from the menu.
.pp
If you choose not to use the default system,
select the Parameter Entry program from the menu (menu item 9).
A "mini" menu lists the three parts of the Parameter Entry program.
Answer the requests for part one, part two, and part three, hitting
ESCAPE when you finish answering the requests for each part. Press
ESCAPE at the "mini" menu to end the Parameter Entry program.  The
Payroll System startup sequence will call up the Account Entry program
automatically, allow you to enter or change account information, then
call up the System Deductions Entry program, and allow you to enter or
change that information.  After the program goes on to Create Sort Parameters,
the system menu returns.  Any further parameter changes can be made
by calling the Parameter Entry, Account Entry, or System Deduction Entry
programs directly from the menu.
.pp
The list below shows selected parameter requests and restrictions on
their responses.  See Chapter 18 for instructions on how to enter numbers
and dates.
.sp
.li
     CO NAME: 60 characters maximum
     ADDRESS: 3 lines of 30 characters each
     CITY: 18 characters maximum
     STATE: Official two-character state abbreviation
     ZIP: 5 digit number only
     GOVT ID NOS.: 10 digit number maximum
     VACATION RATE: 99,999.99 maximum entry
     S.L. RATE: 99,999.99 maximum entry
     RATE1,2,3, and 4 DESCRIPTIVE NAME: 15 characters each
     RATE1 DEFAULT: 99,999.99 maximum entry
     RATE2 FACTOR: 99,999.99 maximum entry
     RATE3 FACTOR: 99,999.99 maximum entry
     MAX PAY FACTOR: 99,999.99 maximum entry
     NORMAL UNITS: 99,999.99 maximum entry
     USER DEFINED PROGRAM NAME: Valid 8 character CP/M filename;
       no asterisks, periods, colons, double quotes, or question
       marks. It may not start with a number.
     USER DEFINED PROGRAM DESC: 30 characters maximum
     CHECK WRITING MAX AMOUNT (BEFORE VOIDING): 9,999,999.99 max
     OFFSET CASH ACCOUNT: 1 to 3 digit Payroll System Account
       Reference Number
     DEFAULT DIST ACCOUNT: 1 to 3 digit Payroll System Account
       Reference Number
.sp2
.ce
DISK DRIVE ALLOCATION
.pp
The default disk drive allocation scheme is presented below (Figure 20.1).
Most
people will not need or want to alter it.  Those with more than
two drives may want to assign the Employee Master file, the Employee
History file, or the Pay Transaction (Hours) file to drive C or D if
they have a large number of employees, or expect many transactions.
.sp
.cp5
.li
  TAX TABLES	  [B]	PAYROLL ACCOUNTS  [B]	TRANS BATCH [B]
  SYSTEM DED/EARN [B]	EMPLOYEE MASTER   [B]	PAY CHECK   [B]
  COMPANY HISTORY [B]	EMPLOYEE HISTORY  [B]	PAY DETAIL  [B]

	(Figure 20.1: Default Disk Drive Allocation Scheme)
.pp
The system creates the two parameter files (PYPR1 and PYPR2)
on the drive you specify
as the System File Drive at first time startup.  It always expects them
on that drive and issues an error message if they are not.  The
System File Drive can be either drive A or B only.  If you chose
the default Payroll System no request is issued for the System File Drive
since drive B is assumed.
Answering the System File Drive request with RETURN defaults to the
currently logged drive. The Sort Parameter file is also created on the
System File Drive.
.pp
If you choose the General Ledger interface option,
never make the GL DATA DRIVE the same as the PAYROLL ACCOUNT FILE.
.pp
The CRT.DEF file, if used, should always reside on the currently logged
drive.
.he
.bp
.ig
CHAPTER 21
T H E	 A C C O U N T	  F I L E
A N D	 G L	I N T E R F A C E
.sp2
.pp5
At first time startup, after parameter entry is complete, the Payroll
System goes automatically to the Account Entry program.  If you
have chosen the default Payroll System, no requests are issued and
the default account file shown below (Figure 21.1) is created.	If you
do not choose the default Payroll System, the Account Entry program gives
you the opportunity to add to or modify the existing account information.
Whether you have chosen the default Payroll System or not, after the
startup sequence is over, you can examine or modify the Payroll Account
file by choosing the Account Entry program from the menu.
.he Operations: Chapter 21
.sp
.li
    Payroll System			       (General Ledger)
 Account Reference No.	    Account Name	Account Number

	1	       CASH			    011100
	2	       ACCRUED LIABILITY	    121100
	3	       SALARY EXPENSE		    411100
	4	       ACCRUED PAYROLL COSTS LIAB.  131100

	    (Figure 21.1: The Default Account File)
.pp
If you have selected the General Ledger interface option the Account
Entry program prompts you to insert the disk with your GL Chart of
Accounts file into the drive you named as the GL DATA DRIVE during
parameter entry.  On a two drive system with the Account file on B,
this would have to be drive A.	Remove the A disk when instructed and
insert the GL data disk.  At the end of Account Entry, switch them
back to the original configuration when instructed.
.pp
When the cursor is waiting at the RECORD NUMBER field, you may enter
the Payroll System Account Reference Number (the relative record number,
see Chapter 9)
of the account you wish to examine or change, or issue one of the
following commands:
.sp
.cp7
.li
	    CTRL-N: Get NEXT account on the file.
	    CTRL-B: Go BACK to previous record.
	    CTRL-A: ADD new accounts to the file.
	    CTRL-S: SAVE newly entered or changed account
		    information to disk.
	    CTRL-R: Refresh garbled DMS screen
	    RETURN: Move cursor to ACCOUNT NUMBER field.
	    ESCAPE: End the Account Entry program.
.sp
Once established, an account may not be deleted from the file, and
should not be reused in the current year.  The CTRL-A command puts
the program in an Add-Records mode.  While in this mode,
after entering the Account Number and Name, hit ESCAPE to add the record
to the file and begin the next.  If no more additions are needed, hit
ESCAPE again to leave the Add-Records mode.  You can then issue any of the
commands above, or ESCAPE once more to end the Account Entry program.
.pp
The CTRL-S command saves additions and modifications to the Account file
out to disk.  Although the Account Entry program saves new account information
out to disk automatically when it ends, the CTRL-S command minimizes the
loss of data (and thus the loss of time and labor) that might occur should
an equipment or other failure occur before the program ends normally.
.pp
The CTRL-R command can be issued from any Display Management System
screen to refresh the screen.  No data is changed and the cursor remains
where it was when the command was given.  You will need this command
rarely, if ever.
.pp
When the cursor is at either the ACCOUNT NUMBER or ACCOUNT NAME field,
the following commands are available:
.sp
.li
	    RETURN: Move cursor to next field.
	    CTRL-B: Move cursor BACK one field.
	    CTRL-X: Cancel current account entry transaction.
	    ESCAPE: Enters the account information and returns
		    the cursor to the Record Number field;
		    when in add-records mode, begins next record
		    or leaves add-records mode.
	    CTRL-R: Refresh garbled DMS screen.
.pp
If you have chosen the General Ledger interface option, the Account
Entry program checks each account number you enter to make sure it
is a valid GL account and displays
the current GL account name in the
Payroll Account Name field as the default.  The fourth digit from the
left should not be a zero (e.g., 411000), since this indicates the
account is a title account and not a balance holding account (see your
GL Reference Manual).
Changing the account name
does not alter your GL chart of accounts file.
.pp
Whether GL interface is selected or not the following rules apply
to account numbers:
.sp
.in5
.ti-3
1. The first digit of a CASH account must be 0 (zero).
.sp
.ti-3
2. The first digit of a LIABILITY account must be 1 (one).
.sp
.ti-3
3. The first digit of an EXPENSE account must be 4 (four).
.qi
.pp
If GL interface has been selected, the Payroll Application program
creates a General Ledger batch file that can be processed by
Applewood Computers's General Ledger after it is PIP'ed over to the GL data disk.
There is no provision for switching disks to receive the GL batch
file as it is output by the Payroll Application program.  If you have
a three or four drive system, however, you can assign the GL batch
file (through Parameter Entry) to a drive separate from the other
Payroll system files.  When you run the Payroll Application program,
it places the GL batch on the disk in that drive.
.he
.bp
.ig
CHAPTER 22
S Y S T E M - W I D E
E A R N I N G S   A N D   D E D U C T I O N S
.sp2
.pp5
To enter or change standard withholding and deduction information,
alter the Federal tax table, and enter or change the five user-definable
earnings and deductions categories, select the menu choice "System
Deduction Entry".  A "mini" menu will appear from which you can make the
desired selection:
.he Operations: Chapter 22
.sp
.li
	1. STANDARD DEDUCTION RATES
	2. FEDERAL WITHHOLDING TAX TABLE ENTRY
	3. SYSTEM EARNING AND DEDUCTION INFORMATION
.pp
After the appropriate Display Management System screen appears, you can
move the cursor along a fixed path from field to field indefinitely,
or until you have made all the entries you wish to make.  RETURN moves
the cursor to the next field, CTRL-B moves the cursor back one field.
ESCAPE at any point returns you to the "mini" menu shown above; ESCAPE
from there returns you to the main system menu.
.sp2
.ce
1. STANDARD DEDUCTION RATES
.pp
Each deduction category can be enabled or disabled systemwide by answering
Yes or No
to the USED? request.  Enter Y for Yes
and N for No (in upper or lower case).
.pp
An account number may also be entered for each
category.  Enter the Payroll System account reference number,
not the General Ledger
account number.
A printout
of the accounts on the account file is available at any time from the
system menu.
.pp
Rate and Limit information for FICA, COMPANY FICA, SDI, COMPANY FUTA,
COMPANY SUI, AND EIC (Earned Income Credit) can be entered with up
to two digits to the right of the decimal.  Commas
are inserted automatically.
.pp
If any employees have requested advance payment of their Earned Income
Credit by filing form W-5 (Earned Income Credit Advance Payment
Certificate), enter the rate used to figure the payment, and the
highest ANNUAL wage upon which the payment may be based.  For example,
if the current rate is 10% of wages up to an annual equivalent wage
of $5,000, enter 10% and $5,000.  The advanced EIC
payment for
employees whose annual equivalent
wage is over a certain limit must be reduced by a given percentage
of the wages over that limit.  (The "annual equivalent wage" is
the amount an employee would earn yearly based on the wage
he or she receives in one pay period.
For example, if an employee
earns $1,000 per month, the annual equivalent wage is $12,000.)
Enter the rate at which the payment is reduced, and the annual equivalent
wage beyond which the reduction must occur.  For example, if the
advanced payment must be reduced by 12.5% of the amount over $6,000,
enter those figures as the EIC EXCESS RATE and LIMIT, respectively.
.sp2
.ce
2. FEDERAL WITHHOLDING TAX TABLE ENTRY
.pp
Federal withholding is calculated by the percentage method, as explained
in the IRS Employer's Tax Guide (Circular E), using the Annual Payroll
Period table.  Applewood Computers distributes the Payroll System with the current
withholding table, but if the table changes you must enter the new
information.  There is a separate table for married and single employees.
Do not enter commas.  The table published by the IRS should appear somewhat
like the one below:
.sp
.cp5
.li
    Over--   But not over--		      of  excess over--
    $1,420   --$3,300	   15%		      --$1,420
    $3,300   --$6,800	   $282.00 plus 18%   --$3,300
    $6,800   --$10,200	   $912.00 plus 21%   --$6,800
    etc...
.sp
Take the figure in the "Over--" column, and enter it into the "Wages
Over" column on the Federal Withholding Tax Table Entry screen.  Take
the corresponding percentage (15%, 18%, 21%, etc.) from the IRS table,
and enter it into the "Percent" column on the Tax Table Entry screen.
For example,
.sp
.cp6
.li
		     S I N G L E
	   WAGES OVER		 PERCENT
	 [	 1,420]      [		15]
	 [	 3,300]      [		18]
	 [	 6,800]      [		21]
	 etc...
.sp2
.ce
3. SYSTEM EARNING AND DEDUCTION INFORMATION
.pp
To create a system-wide earnings or deduction category, enter a 15
character description, and state whether it is an earning or deduction
(E=Earning, D=Deduction).  For a deduction,
enter the number of the Payroll System
account to which the liablity will be credited. For earnings and deductions
enter the
Check Category where it should be printed (see Chapters 15 and 28), and
enter the information required to calculate the earning or deduction.
.pp
If there is a limit to the amount that can be earned or deducted in
one year, enter a Y for Yes to the LIMIT? request, and enter the amount of the
limit into the LIMIT field.  If you answer the Limit Used? (LIMIT?) request
with N (No),
the earning or deduction is calculated every pay period.
.pp
The FACTOR can be either a flat amount to be added or subtracted
each pay period, or a percentage of the employee's gross pay.
If a flat amount, enter A (Amount)
at the A/P? request; if a percentage, enter a P (Percentage).  The factor
may contain two digits to the right of the decimal. Do not
enter commas.
.pp
To enter the tax status of an earnings category, key one of the numbers
below into the EXCLD (Excluded) field:
.sp
.cp7
.li
	1   Fully taxable.
	2   Subject to all taxes except FICA.
	3   Subject to all taxes except Federal,
	    state, and local income taxes.
	4   Not subject to taxation.

	 (Figure 22.1: Tax Exclusion Code)
.sp2
.in8
.ti-8
EXAMPLE: If a $500 bonus is to be paid to all employees at the rate
of $50 per pay period
you would make the entry below.  No account number is needed since
earnings are distributed according to the distribution accounts
set for each employee.
.qi
.ll72
.sp
.cp4
.li
USED? E/D DESC	   E/D	ACCT  A/P    FACTOR LIMITED  LIMIT XLCD CAT

 [Y][BONUS	 ] [E] [   ]  [A] [    50.00] [Y] [  500.00][1][ 6]
.ll65
.sp2
.in8
.ti-8
EXAMPLE: If 2% of the employees' gross pay is to be deducted for a
supplemental health care plan every pay period, make the entry
as shown here:
.qi
.sp
.ll72
.li
[Y][SUPP HEALTH ] [D] [  2]  [P] [     2.00] [N] [	  ][  ][15]
.ll65
.qi
.he
.po8
.bp
.ig
CHAPTER 23
E N T E R I N G   A N D    C H A N G I N G
E M P L O Y E E    I N F O R M A T I O N
.sp2
.pp5
Putting an employee on the Payroll System is a three step process:
.he Operations: Chapter 23
.sp
.li
	1. Create an employee record using
	   the EMPLOYEE ENTRY program.

	2. Enter deduction/earnings and cost
	   information using the EMPLOYEE
	   DED/EARN AND COST program.

	3. Enter pay rate, vacation rate, sick
	   leave rate, and other information
	   with the EMPLOYEE RATE ENTRY program.
.sp
Steps one and two are covered in this chapter, step three in the following
chapter (Chapter 24). The information entered through these three programs
is stored on the Employee Master file.
.sp2
.ce
EMPLOYEE ENTRY
.pp5
When you select the Employee Entry program from the Payroll system
menu the Employee Entry screen (see Chapter 32) appears with the cursor
waiting at the EMPLOYEE NUMBER field.
From there
you can give commands to change, examine,
or save employee records:
.sp
.li
	  CTRL-A:  ADD records to the employee file.
	  CTRL-N:  Get the NEXT record on the file.
	  CTRL-B:  Go BACK to the previous record.
	  CTRL-S:  SAVE newly entered or changed records
		   to disk.
	  CTRL-R:  REFRESH garbled DMS screen.
	  ESCAPE:  End the Employee Entry program.
.sp2
.ce
ADD AN EMPLOYEE RECORD
.pp
To add an employee to the file, give the command
CTRL-A when the cursor is at the EMPLOYEE NUMBER field.
The program finds the first available employee number, displays it in
the Employee Number field, and moves the cursor to the NAME field
where you can begin entering data for the employee.
.pp
Employee
numbers are assigned automatically; the lowest available number is always
assigned first.  Employee numbers take the following form:
.sp
.li
		  nnnnc
.sp
where "nnnn" is any number from 1 to 9999 and "c" is a "check digit"
that protects against miskeys and transposition errors.  When
an employee number is requested, enter all five digits.  The relative
position of the record on the file can be easily determined by
ignoring the check digit.
.pp
Press ESCAPE after entering the data for each employee.  The program
will display the next available employee number and mode the cursor
to the Name field where you can begin the next record.
An ESCAPE at the first request of the
next record takes the program out of the Add-Records mode and returns
the cursor to the Employee Number field.  ESCAPE from there ends the
Employee Entry program.
.in6
.ti-6
.sp
NAME: Enter the name in the form:
.sp
.li
		   Firstname Middlename Lastname
.in6
Up to 30 characters are allowed.
.ti-6
.sp
ADDRESS: Up to 30 characters per line.
.ti-6
.sp
CITY: 18 characters maximum.
.ti-6
.sp
STATE: Official two character abbreviation.  If the state is California (CA),
special requests are issued for head of household status and number of special
withholding allowances.
.ti-6
.sp
ZIP: Five digits.  This field accepts only numbers.
.ti-6
.sp
SOCIAL SECURITY NO.: Accepts only the numbers 1 through 9, and dashes
(-). Must have the format (nnn-nn-nnnn).
.ti-6
.sp
BANK ACCOUNT NO.: Accepts only the numbers 1 through 9, and dashes (-).
Optional entry.  25 characters maximum.
.ti-6
.sp
BIRTH DATE: Accepts only the numbers 1 through 9, and slashes (/).
Enter in the form MM/DD/YY or DD/MM/YY, depending on the date
format you chose in Parameter Entry.
.ti-6
.sp
SEX: Male = M, Female = F.
.ti-6
.sp
PENSION (Y/N): Y = Yes, the employee is on a company pension plan;
N = No, the employee is not on a company pension plan.
.ti-6
.sp
JOB CODE: 3 characters, letters or numbers. Informational only.
.ti-6
.sp
START DATE: Enter MM/DD/YY or DD/MM/YY, depending on the date format
you chose in Parameter Entry.
.ti-6
.sp
TERM. DATE: Enter MM/DD/YY or DD/MM/YY, depending on the date format
you chose in the Parameter Entry.
.ti-6
.sp
EMPLOYMENT STATUS: A = Active, L = On Leave, T = Terminated.
Transactions may not be entered for employess who are On Leave
or Terminated (unless the Display Employee Name option in Transaction
Entry is declined). Autogen pay is suppressed for Employees On Leave.
A terminated employee's record remains on the file until year's end.
After you print a (W-2) form , and Close for the Year the employee's
number is available for reassignment.
.qi
.sp2
.ce
CHANGE AN EMPLOYEE RECORD
.pp
To the change information on an employee record,
with the cursor at the Employee Number field call up the record for the
employee by any combination of the following methods:
.sp
.li
	     1. Enter the 5 digit Employee Number.
	     2. Give the CTRL-N, get next record command.
	     3. Give the CTRL-B, get previous record command.
.sp
If no record is present,
CTRL-N gets the first record on the file, CTRL-B the last.
Hit RETURN to move the cursor to the field you want to change and type
in the new response.
Outside of the Employee Number field CTRL-B backs the cursor
up one field.  ESCAPE at
any point enters the changes onto the file and returns the
cursor to the Employee Number field where you can call up more records or
issue other commands.
.sp2
.ce
EXAMINE AN EMPLOYEE RECORD
.pp
Call up the record as explained above
for changing a record.
.sp2
.ce
SAVE NEWLY ENTERED OR CHANGED RECORDS
.pp
The SAVE command, which can
be given when the cursor is at the Employee Number field, updates the
disk file with the newly created records or the changes made to the
old records.  Diligent use of this command can save
you hours of work should a system failure occur due to a hardware
or other type of problem (such as a pulled plug or write protected
disk).	As a rule of thumb, use the Save command as soon as you
enter more records than you would care to enter again. One hard
lesson is usually enough to demonstrate the value of this command.
.cp12
.sp3
.ce
EMPLOYEE DED/EARN AND COST DIST.
.pp
The Employee Deduction/Earnings and Cost Dist. program is used
to specify the accounts to which an employee's gross earnings and
associated costs should be distributed, to exempt an employee from
system-wide earnings or deductions, or to enter earnings or deductions
that apply to a particular employee only.
.pp
When you select the program from the menu an entry screen appears with the
cursor waiting at the EMPLOYEE NUMBER field.  You may call up an
employee record by any combination of the following methods:
.sp
.li
		1. Enter the five digit employee number.

		2. Give the CTRL-N command to call up the
		   first record on the file, or if a record
		   is already present, the next record.

		3. Give the CTRL-B command to call up the
		   last record on the file, or if a record
		   is already present, the previous record.
.in6
.sp
.ti-6
GL COST DISTRIBUTION OF EMPLOYEE GROSS AND EMPLOYER COST: The total
cost of an employee to the employer can be distributed to up to five ledger
accounts.  Enter the Payroll System Account Reference Number and the percentage
of the total to be distributed to that account.  The total of the percentages
must equal 100.00 before you can go on to the next employee.  For example,
if the employee cost is to be divided evenly between three accounts,
the following entry would be correct:
.sp
.li
	     ACCOUNT [ 10]    PERCENT  [   33.33]
	     ACCOUNT [ 11]    PERCENT  [   33.33]
	     ACCOUNT [ 12]    PERCENT  [   33.34]
	     ACCOUNT [	 ]    PERCENT  [	]
	     ACCOUNT [	 ]    PERCENT  [	]
			      TOTAL    [  100.00]
.sp
.ti-6
EXEMPTIONS FROM TAX WITHHOLDING:  To exempt an individual employee
from Federal, state, local, FICA, or SDI withholding, enter Y for
Yes, the employee is exempt, and N for No, the employee is not exempt.
.sp
.ti-6
EXEMPTIONS FROM SYSTEM DEDUCTIONS: To exempt an individual employee from
any of the five user-defined system-wide earnings or deduction
categories, enter
a Y for Yes, the employee is exempt, or N for No, the employee is not
exempt.
.sp
.ti-6
EMPLOYEE SPECIFIC DEDUCTIONS: Up to three earnings
or deduction categories can be established for a specific employee.
The procedure is exactly the same as described for creating the systemwide
earnings or deductions.  See Chapter 23 for
instructions.
.qi
.pp
After entering or adjusting the information in the entry fields,
an ESCAPE will update the employee file and return the cursor to
the Employee Number field.  From there, you can call up another employee
record, or press ESCAPE to end the program.
.he
.bp
.ig
CHAPTER 24
E N T E R I N G    A N D    C H A N G I N G
P A Y	R A T E    I N F O R M A T I O N
.sp2
.pp5
When you select the Employee Rate Entry program from the system menu,
the Employee Rate Entry screen (see Chapter 32) appears with the cursor
waiting at the Employee Number field.  You can call up the record for an
employee by any combination of the following methods:
.he Operations: Chapter 24
.sp
.li
	   1. Enter the five digit employee number.
	   2. Give the CTRL-N command to call up the
	      first record on the file, or if a record
	      is already present, the next record.
	   3. Give the CTRL-B command to call up the
	      last record on the file, or if a record
	      is already present, the previous record.
.sp
After the program displays the record, hit RETURN to begin entering data.
After answering the requests for an employee, press ESCAPE to return the
cursor to the Employee Number field.  You can then call up additional
records or press ESCAPE again to end the program.
.sp2
.ce
PAY INTERVAL AND PAY RATES
.in6
.sp
.ti-6
PAY TYPE: H = Hourly, S = Salaried.
.sp
.ti-6
PAY INTERVAL: Enter the frequency with which the employee is paid:
.sp
.li
	       W = Weekly
	       B = Bi-weekly
	       M = Monthly
	       S = Semi-monthly
.sp
.ti-6
RATE1 THROUGH RATE4: The descriptive names of Rate1 through Rate4
set in Parameter Entry
appear beside the pay rates currently in effect.  If the default or
current value is acceptable, you can hit RETURN at each field.
.pp
The largest values allowed in the Rate fields are as follows:
.sp
.li
	       Rate1.........99,999.99
	       Rate2.........99,999.99
	       Rate3.........99,999.99
	       Rate4.........99,999.99
.sp
Do not enter commas.
.sp
.ti-6
RATE4 TAX EXCLUSION TYPE:  See Chapter 36, Command Summary for the
Tax Exclusion Type code.
.sp
.ti-6
AUTOGEN UNITS: Enter the number of units per pay period.  This is
usually 1 (one) for salaried employees.  If an hourly employee is to
receive automatically generated pay enter the number of hours per pay
period (for example, 40 for weekly pay, 80 for bi-weekly, etc.).
If the employee is not to receive automatic pay, enter 0 (zero) autogen
units.
.sp
.ti-6
NORMAL UNITS: Enter the number of Rate1 pay units an employee
normally receives in one pay period.  When a Rate0 transaction is entered
through the Pay Transaction entry program, the employee will automatically
be paid at Rate2 for all units over the normal units value.
The largest value allowed in the
Normal Units field is 999.99.
.sp
MAXIMUM PAY: Warnings appear on screen during payroll application and
journal printout for
employees whose gross pay exceeds this amount in one
pay period.
The largest value allowed in this field is 999,999.99.
.qi
.sp2
.ce
TAX WITHHOLDING ALLOWANCES
.in6
.sp
.ti-6
MARITAL STATUS: M = Married, S = Single.
.sp
.ti-6
FEDERAL WH ALLOWANCES:
.ti-6
STATE	WH ALLOWANCES:
.ti-6
LOCAL	WH ALLOWANCES: Enter the number of withholding allowances claimed
by the employee for Federal, state, and local tax purposes.  These fields
accept only numbers from 1 to 99.
.qi
.sp2
.ce
STATE OF CALIFORNIA SPECIAL INFORMATION
.in6
.sp
.ti-6
HEAD OF HOUSEHOLD: Enter Y for Yes, the employee is considered the head
of household for California tax purposes,
or N for No, the employee is not considered
the head of household.
.sp
.ti-6
SPECIAL WITHHOLDING ALLOWANCES: Enter the number of special withholding
allowances claimed by the employee for California tax purposes.  This
field accepts only numbers from 1 to 99.
.qi
.sp2
.ce
VACATION, SICK LEAVE, AND COMP TIME INFORMATION
.in6
.sp
.ti-6
VACATION RATE: Enter the number of vacation units the employee should
receive every pay period.  For example, if an employee earns 10 days
of vacation per year, and is paid on a monthly schedule, enter the
vacation rate as 10/12 = .83 days per month.  If an employee earns
80 hours of vacation per year and is paid bi-weekly, calculate the
number of hours of vacation he or she earns every two weeks and enter
that number as the vacation rate (80/26 = 3.08).
The largest value allowed by this field is 99,999.99.  Only two digits
to the right of the decimal are allowed.
.sp
.ti-6
VACATION ACCUMULATED:  Enter the number of vacation units the employee
has accumulated. Note that this is not the number of units the
employee has remaining.
The largest value accepted by the field is 99,999.99.
.sp
.ti-6
VACATION USED:	Enter the number of vacation units the employee has
used so far this year.	The difference between Vacation Used and
Vacation Accumulated is the number of vacation units the
employee has remaining.
.sp
.ti-6
SICK LEAVE RATE: Enter the sick leave rate as explained above for
the Vaction Rate.  The largest value accepted by the field is
99,999.99.
.sp
.ti-6
SICK LEAVE ACCUMULATED: Enter the number of sick leave units the employee
has accumulated,
as explained above for Vacation Accumulated.  The largest value accepted
by the field is 99,999.99.
.sp
.ti-6
SICK LEAVE USED: Enter the number of sick leave units the employee
has used so far this year as explained above for Vacation Used.
The difference between Used and Accumulated is the number of sick leave
units the employee has remaining.
.sp
.ti-6
COMP TIME ACCUMULATED: Enter the
number of comp time units he or she has accumulated so far this year.
This is not the number of comp time units the employee has remaining.
The largest value accepted by the field
is 99,999.99.
.sp
.ti-6
COMP TIME USED:  Enter the amount of comp time the employee
has used so far this year.
The difference between Comp Time Accumulated and Comp Time Used
is the amount of compensatory time the employee has remaining.
The largest value accepted by this field is 99,999.99.
.qi
.he
.bp
.ig
CHAPTER 25
T H E	P A Y R O L L	P R O C E S S I N G    C Y C L E
.sp2
.pp5
Once the setup steps have been completed you can begin the regular
processing cycle you will repeat every pay period.  There are seven
steps in the payroll processing cycle:
.he Operations: Chapter 25
.sp
.li
	   1. Enter Pay Transactions
	   2. Print a Pay Transaction Batch Proof
	   3. Back up all files.
	   4. Sort the Pay Transaction batch.
	   5. Run the Apply Payroll program.
	   6. Print the Payroll Journal.
	   7. Print Paychecks and/or the Paycheck Register.
.sp
Chapter 26 covers the first four steps.  Chapter 27 handles
five and six, and Chapter 28 step seven.
.sp
.in6
.ti-4
(1) ENTER PAY TRANSACTIONS: To do this step, choose the Enter Pay
Transactions program from the system menu.  This program creates a
"batch" file of pay transactions; the file is named PYHRS.001 when
unsorted and PYHRS.101 after it has been sorted. You can add transactions
to the
batch, change transactions, delete transactions, and examine transactions.
Working proofs of the batch may be printed at any time until the final audit
proof is taken.  The file is completely consumed by the
Apply Payroll program.
.sp
.ti-4
(2) PRINT A PAY TRANSACTION AUDIT PROOF:  When the Pay Transaction
batch is accurate and complete, run the Pay Transaction Proof program
from the system menu to generate a hardcopy audit proof.
.sp
.ti-4
(3) BACK UP ALL FILES:	Before taking any further steps, copy all of your
disks.	Use the full disk copy program provided with the Payroll System
(SDCOPY) for copying single density 8" IBM standard format disks, or
some other copy program such as PIP.
.sp
.ti-4
(4) SORT THE PAY TRANSACTION BATCH:  To sort the batch of pay transactions,
choose the Sort Pay Transactions program from the system menu.	Be sure
there is room on the disk with the Pay Transaction file for both the
sorted and unsorted version, and for temporary workfiles whose total size
will never be larger than
the input file.  Use the CP/M STAT command to find the
size of the input (unsorted) PYHRS.001 file, and the amount of space
remaining on the disk (for example, STAT PYHRS.001).
If your CP/M derviative operating system does not
have the SUBMIT command, run the sort outside the Payroll System as
explained in Chapter 26.
.sp
.ti-4
(5) RUN THE PAYROLL APPLICATION PROGRAM:  Call the Apply Payroll program
from the system menu, enter the date you wish printed on the checks, and
state whether the application is for month-based or week-based employees.
Make sure there is room on the disk for both the old and the new Pay Detail
files, since both exist at the same time during processing. The program may
take some time to run with many employees on the system.
.sp
.ti-4
(6) PRINT THE PAYROLL JOURNAL: Before the next payroll application can
take place you must print at least one copy of the Payroll Journal.
Choose the Print Payroll Journal program from the system menu.
.sp
.ti-4
(7) PRINT PAYCHECKS AND/OR THE PAYCHECK REGISTER: If you have chosen
automatic paycheck printing, load the printer with the paychecks, print
a dummy paycheck until they are aligned properly, and then print the
paychecks for employees to be paid that period.  Also print the
paycheck register.  If you did not choose automatic paycheck printing,
print the Paycheck Register so you can prepare checks
from it.
.qi
.he
.bp
.ig
CHAPTER 26
E N T E R I N G    E M P L O Y E E    P A Y
T R A N S A C T I O N S
.sp2
.pp5
When you select the Pay Transaction Entry program from the Payroll
System menu, if no transaction batch exists the program asks
if you want to create one.  Answer Yes to create a new batch, and
No if for some reason (such as if the wrong disk is inserted)
you do not.  If you answer No, the Pay Transaction Entry program ends
and the menu returns. If you answer Yes, the program asks if
you want to Display Employee Names.  If you answer Yes, the
program displays an employee's name beside his or her  employee number
to help make sure you enter the correct transaction information
for the employee.  The Employee Name option also causes the Pay
Transaction Entry program to reject transactions
for Terminated employees, and issuea a warning when a transaction is
entered for an employee On Leave.
.he Operations: Chapter 26
.pp
After you answer the employee name option request, the transaction
entry screen appears with the cursor waiting in the Record Number
field.	The following commands are available while the cursor
is in the Record Number field:
.sp
.in11
.ti-8
CTRL-A: ADD transactions to the file.
.ti-8
CTRL-N: Gets the NEXT record on the file; when no records have been
accessed, gets the first transaction record.
.ti-8
CTRL-B: Goes BACK one record on the file; when no records have been
accessed, brings up the last transaction on the file.
.ti-8
CTRL-D: DELETE the transaction.
.ti-8
CTRL-S: SAVE newly entered or changed records to disk.
.ti-8
CTRL-R: REFRESH garbled DMS screen.
.ti-8
RETURN: Move the cursor to the first request
(when a record has been accessed).
.ti-8
ESCAPE: End the Pay Transaction Entry program.
.ti-8
ENTER RELATIVE RECORD NUMBER: To call up a particular transaction,
enter the position of the transaction record relative to the beginning
of the file. Transactions are stored in the order they are entered.
For example, the third record entered is third on the file, and thus
has the relative record number of 3.
.qi
.sp2
.cp6
.ce
ADDING TRANSACTIONS
.pp
To add transactions to the file give the command CTRL-A.  This puts
the program in the Add-Records mode and moves the cursor to the
Employee Number field.	Enter the 5 digit employee number, the
transaction code, and the number of units (dollars, hours, pieces, etc.),
and a date.  Figure 26.1 lists the Payroll System transaction codes and
the check category where the item prints, if applicable.
Refer to Chapter 13 for a description of each transaction type.
The date can be the date you enter the transaction,
or any reference date you choose.  The default date is the run date
you give when you call up the Payroll System.
.pp
While in the Add-Records mode, RETURN moves the cursor forward one field,
and CTRL-B moves it back one field.  After you answer all the requests for
a transaction, you can move through the fields indefinitely with RETURN
or CTRL-B to adjust the information, or you can press ESCAPE at any point
to begin the next record.  ESCAPE enters the transaction onto the file
and returns the cursor to the Employee Number field, where you can enter
another employee number and begin a new transaction.  To leave the Add-Records
mode, after hitting ESCAPE to enter the transaction, press ESCAPE at the first
request of the next transaction.
This moves the cursor to the Record Number field where you can give one of
the commands listed above, or ESCAPE again to end the program.
.sp2
.cp30
.li
	PAYROLL SYSTEM TRANSACTION CODES

    TRANSACTION     TRANSACTION        CHECK
       CODE	       TYPE	      CATEGORY

	(0)	  Rate0 Pay
	(1)	  Rate1 Pay		 2
	(2)	  Rate2 Pay		 3
	(3)	  Rate3 Pay		 4
	(4)	  Rate4 Pay		 5
	(5)	  Vacation Taken
	(6)	  Sick Leave Taken
	(7)	  Comp Time Taken
	(8)	  Comp Time Earned
	(9)	  Bonus 		 6
       (10)	  Tips Collected	 7
       (11)	  Advance		 6
       (12)	  Sick Pay		 6
       (13)	  Vacation Pay		 6
       (14)	  Other Excluded Pay	 6
       (15)	  Expense Reimbursement  6
       (17)	  Other Pay		 6
       (18)	  Tips Reported 	14
       (27)	  Advance Repay 	14
       (28)	  FWT Add-On		 9
       (29)	  SWT Add-On		10
       (30)	  LWT Add-On		11
       (31)	  FICA Add-On		12

  (Figure 26.1 Transaction Codes and Check Categories)
.sp2
.ce
CHANGING TRANSACTIONS
.pp
To change a transaction entered previously, first locate it by giving
the CTRL-N,
CTRL-B, or Relative Record Number command at the Record Number field.
When the transaction you wish to change is showing, hit RETURN to move
the cursor from the Record Number field to the Employee Number field.
Press RETURN to move the cursor forward to the response you wish to change,
and type in the new response.  You can use the CTRL-B command to move
the cusor back one field.  When the changes are complete, hit ESCAPE
to return the cursor to the Record Number field where you can issue one
of the commands listed above.
.sp2
.cp6
.ce
EXAMINING TRANSACTIONS
.pp
To examine previously entered transactions you simply locate the
transaction by giving one of the following commands from the the
Record Number field:
.sp
.in6
.ti-3
1. Enter the Relative Record Number of the Transaction. For example,
to bring up the fifteenth record on the file, you would enter 15
and hit RETURN.
.sp
.ti-3
2. Give the CTRL-N command to display the NEXT transaction on the
file; gets the first record if none are showing.
.sp
.ti-3
3. Give the CTRL-B command to display the previous transaction on
the file ("Go BACK one record"); gets the last record is none are
showing.
.qi
.sp2
.cp6
.ce
DELETING TRANSACTIONS
.pp
To delete a transaction, first bring up the transaction by the
CTRL-N, CTRL-B, or Relative Record Number command, and then
give the command,
.sp
.li
		CTRL-D
.sp
Once you delete a transaction, you can no longer access it.
The remaining transactions keep their same relative record numbers
since the deleted transaction is not physically removed from the file,
but only marked as deleted.
.sp2
.cp6
.ce
SAVING NEWLY ENTERED OR CHANGED TRANSACTIONS
.pp
The CTRL-S command protects you against the possible loss of perhaps hours
of labor in the event of an equipment failure.	The CTRL-S command
saves the transactions added or changed during the current
run of the program out to disk.  Should the program end prematurely
for any reason, all transactions saved up to the point of system
failure would be stored safely on disk.  A rule of thumb is to give
the CTRL-S command when you have keyed in more transactions
than you would care to enter again.  It usually takes one hard lesson
to appreciate the value of this command.
.sp2
.cp4
.ce
REFRESH GARBLED DMS SCREEN
.pp
If for some reason a Display Management System screen should become
garbled, the CTRL-R command refreshes it.  The cursor position
and data showing on the screen remain unchanged.  You can issued
this command at any point.  You will rarely, if ever, need to use it.
.sp2
.cp6
.ce
PROOFING THE TRANSACTION BATCH
.pp
At any point during the development of the Pay Transaction batch file
you can order a numbered printout of the transactions exactly as
they have been entered.  To do so, choose the Transaction Proof
program from the Payroll System menu.  When you finish entering
transactions for all employees to be paid on the next payroll application,
print a final proof report and save it as part of your audit trail.
If you enter additional transactions, the sorting program will warn you
that not all of the transactions have been proofed, and give you an
opportunity to do so.  Check your transaction entries very carefully to
make sure all employees will be paid the correct amounts.  Once a
Payroll Application is run, it cannot be run over, except from your backup
files.
.sp2
.cp6
.ce
SORTING THE TRANSACTION BATCH
.pp
Before it can be processed by the Apply Payroll program, the
Pay Transaction file must be sorted into employee number order.
To do so, select the Sort the Transaction Batch program from the
system menu.  Be sure there is room on the disk for both the sorted
and unsorted versions (they are of equal size), and if the file
is large, for temporary workfiles whose total size will never exceed
the size of the input (unsorted) file.	 Use the CP/M STAT command
to find the size of the input file and the amount of space remaining on
the disk. When the STAT command is on drive A and you wish to find the
space available on B, give the command:
.sp
.li
		A>STAT B:
.sp
and follow with RETURN.  To find the size of the unsorted transaction
batch, you would enter:
.sp
.li
		A>STAT B:PYHRS.001
.sp
since PYHRS.001 is the name of the unsorted file.
.pp
Because the sorting program (QSORT.COM) is not a CBASIC language program,
when you choose the Sort Transaction Batch program from the menu, the
Payroll System ends automatically, performs the sort under the CP/M
environment, and then returns.	To do this it uses the SUBMIT command.
If your CP/M derivative operating system does not have this command (most
do), you
cannot sort the transaction batch from the menu.  Instead, you must
end the Payroll System, and with QSORT on the currently logged drive, and
the Sort Parameter file
on the System File Drive (if B), type:
.sp
.li
		A>QSORT B:PYHOURS
.sp
When sorting is complete, re-enter the Payroll System by typing
CRUN2 PY.
.he
.bp
.ig
CHAPTER 27
P A Y R O L L	 A P P L I C A T I O N
.sp2
.pp5
The Apply Payroll program is the central processing program of
the Payroll System.  It is the program that calculates earnings
and deductions, updates historical information, and creates
a check file to be used by the check printing program.	If General
Ledger interface is selected, it creates the batch of transactions
that can be processed without alteration by the Applewood Computers General Ledger.
.he Operations: Chapter 27
.pp
Before running the Apply Payroll program be certain the following
steps have been taken:
.sp
.in6
.ti-3
1. MAKE SURE THE PAY TRANSACTION BATCH FILE IS COMPLETE AND ACCURATE.
Study the final transaction proof (the audit proof) to be certain
that pay transactions have been entered correctly or an appropriate
number of Autogen Units set for each employee. You cannot run the
Apply Payroll program over again, except by using your backup files.
.sp
.ti-3
2. MAKE SURE A FINAL (AUDIT) PROOF OF THE BATCH HAS BEEN PRINTED.
The proof should be taken before the batch is sorted.
.sp
.ti-3
3. MAKE SURE YOU HAVE MADE DUPLICATE COPIES OF EVERY DISK.  This "backing
up" is crucial to many error recovery procedures. If you forget to
pay an employee or enter incorrect transaction information, you must have
these backup disks to repair the error.  Copy all of your processing disks
after printing the audit proof and before sorting the transaction batch.
See Chapter 35 for error recovery procedures and instructions
for using SDCOPY, the single-density full-disk copying program distributed
with the Payroll System. If you neglect this step error recovery will
consume much more time than it should.	The information on the disk
is far more valuable than the disk itself.
.sp
.ti-3
4. MAKE SURE THE BATCH HAS BEEN SORTED BY EMPLOYEE NUMBER.  See Chapter
26 for instructions.
.qi
.sp2
.cp10
.ce
RUNNING THE APPLY PAYROLL PROGRAM
.pp
To apply payroll for the period, select the Apply Payroll program
from the system menu and enter the Check Date.
This is the pay date that will be printed on the checks, and also
the reference date of the General Ledger transactions.	It is the
date you incur withholding liabilities.  See Chapter 18 for instructions
on how to enter a date.
.pp
If the program cannot tell, it will ask you to state whether the
application is for week-based or month-based employees.  Enter
W for week-based and M for month based.  Type "RUN" to begin the
processing.
.pp
The Apply Payroll program knows whether to process
"fast" interval employees (weekly or semi-monthly), "slow" interval
employees (bi-weekly or monthly), or both fast and slow interval employees
by checking Ending Day of Last Period for the interval.
.he
.bp
.ig
CHAPTER 28
P R I N T I N G    P A Y C H E C K S
A N D	 T H E	  P A Y C H E C K    R E G I S T E R
.sp2
.pp5
To print the paychecks, select the Print Pay Checks program from the
system menu. The program displays the last check number and
requests the starting check number.  This can be any number you want.
.he Operations: Chapter 28
.pp
The program gives you an opportunity to align the check forms in
the printer.  Type Y or hit RETURN to print a test form.  When the
forms are aligned properly, respond to the test form request with N.
.pp
Should your forms jam in the printer, you can press ESCAPE to stop
or RETURN to continue.	There is no facility for reprinting individual
checks.  If a check is spoiled rerun the printout from the
beginning.  If someone didn't get paid or got paid the wrong amount,
rerun the application with your backup files.
.pp
To print the check register, select the Print Pay Check Register program
from the system menu.	Be sure your paper is aligned properly.
.he
.bp
.ig
CHAPTER 29
E N D I N G    T H E	Q U A R T E R
.sp2
.pp5
The end of quarter processing steps should be taken after
the last payroll application of the quarter and before the first
payroll application of the next quarter.  The Apply Payroll program
will not let you continue if the Check Date you enter falls within
the new quarter and the quarter end processing steps have not been
taken.	These steps are:
.he Operations: Chapter 29
.sp
.li
		1. Print 941 Report.
		2. Run Close for the Quarter program.

	     (Figure 29.1: Quarter End Processing Steps)
.sp2
.ce
PRINT 941 REPORT
.pp
To print the Federal 941 report, Employer's Quarterly Federal Tax
Return, select the Print 941 Report program from the system menu.
The first half of a facsimilie reproduction of the 941 report
appears on screen with the current values inserted in the appropriate
fields.  If you wish to use the stick-on label provided with the
government form, you can suppress printing of your company's name and
address information by answering No (N) to the PRINT NAME AND ADDRESS
request.
.pp
You may change two of the fields on this screen:
ADJUSTMENT OF WITHHELD INCOME TAX FROM PREVIOUS QUARTERS and ADJUSTMENT
OF FICA TAXES.	The values in these fields are subtracted from the
INCOME TAX WITHHELD and TOTAL FICA TAXES figures respectively.
Hit RETURN until the cursor is at the field you wish to change and type
in the new value.  The new adjusted totals are computed automatically.
When the data on this screen is complete and accurate, press ESCAPE to
bring up the second half of the facsimilie form.
.pp
This screen lets you change the date of deposit and the amount
deposited to conform to the actual dates and amounts.  The program
assumes your deposits are the same as your liabilies for the period.
If this is not the case, move the cursor forward  to the fields that need to
be changed by RETURN (or CTRL-B to go back), and type in the new value.
The totals are re-computed automatically.  Hit ESCAPE when the data is
correct and follow the instructions to align the form in your printer.
.pp
To align the form, insert it into the printer so that the line of
dashes prints directly on top of the black
line below the space for the company name and address.	If the line
prints high or low, move the form up or down in the printer accordingly
until it is aligned properly.  You may repeat the alignment process
until it is correct.  When the form is aligned properly, press RETURN
as instructed.	Alternatively, print the report on plain paper and
transfer the information  by hand.
You can run the Print 941 Form program as many
times you want.
.sp2
.ce
RUN CLOSE FOR THE QUARTER PROGRAM
.pp
After you print the 941 form, run the Close For the Quarter program.
This program issues no requests, except for confirmation that it should
be run.  The program will not continue if the 941 form has not been
printed, or not enough applications have taken place.  New interals
may only be added at the start of a new quarter.
.he
.bp
.ig
CHAPTER 30
E N D I N G    T H E	Y E A R
.sp2
.pp5
The end of the year processing steps should be taken after end of
quarter processing is complete for the fourth quarter, and before
the first payroll application of the new year.	The Apply Payroll program
will not let you continue if the Check Date you enter falls within the
new year and the year end processing steps have not been taken.
.he Operations: Chapter 30
.sp
.li
		1. Close for the fourth quarter.
		2. Print 940 Report
		3. Print W-2 forms for all employees.
		4. Run Close For the Year program.

	      (Figure 30.1: Year End Processing Steps)
.sp2
.ce
PRINT 940 REPORT
.pp
To print the Federal 940 Report, after
closing the fourth quarter (October, November, December) in the usual
way (see Chapter 29), select the Print 940 Report program from the system
menu.  Since the 940 report is designed to be printed on plain paper and
the information
transferred to the government form manually, check to make sure the
paper is aligned correctly in your printer, and press RETURN to begin
printing.
.sp2
.ce
PRINT W-2 FORMS FOR ALL EMPLOYEES
.pp
To print W-2 forms, select the Print W-2 Report program from the system
menu.  A screen will appear from which you can choose various printing
options:
.sp
.in6
.ti-3
1. Select ALL EMPLOYEES, RANGE OF EMPLOYEES, TERMINATED EMPLOYEES, or
A SINGLE EMPLOYEE.  ALL EMPLOYEES instructs the program to print W-2
forms for all Active, On Leave, and Terminated employees.  The ALL
EMPLOYEES selection must be chosen before you are allowed to close for
the year.
If you
select a RANGE OF EMPLOYEES, you will be asked to enter the starting
employee number and the ending employee number.  The program checks
to make sure the range you enter is valid.  If you select A SINGLE
EMPLOYEE, the program asks for the five digit employee number. Selecting
TERMINATED EMPLOYEES causes the program to print W-2 forms for terminated
employees only. This can be done at any time during the year.
.sp
.ti-3
2. PRINT STATE AND LOCAL AMOUNT: To suppress printing of the
state and local withholding information, enter N for No.
.sp
.ti-3
3. NUMBER OF FORMS PER EMPLOYEE: Enter the number of forms to
print for each employee.  For example, when you print on three
part carbon forms (one part for the Federal return,
one for the state, and one
for the employee's personal records), you would enter 1.  If you
do not use carbon forms, and want to print one form for Federal, one for
state, one for local, and one for the employee, you would enter 4.
If a form for local income tax is not needed you would enter 3.
.sp
.ti-3
4. NUMBER OF LINES BETWEEN FORMS: Enter the number of lines between the
last print line of the previous form and the first print line of the
next form.
.sp
.ti-3
5. STARTING PRINT COLUMN: If you need to shift the
printing to the left or the right, add or subtract the required number
of positions to or from the current value.
.qi
.pp
After answering the above requests, align your continuous form W-2's
in the printer and answer Y for Yes to the PRINT TEST FORM question.
When you hit RETURN a dummy W-2 form will be printed.  Adjust the alignment
of the forms in the printer and print another dummy form.  If the alignment
is correct, answer N for No to the PRINT TEST FORM request and hit RETURN.
Otherwise repeat the alignment process until the forms are aligned
properly.
.sp2
.ce
RUN CLOSE FOR THE YEAR PROGRAM
.pp
After you print the 940 form and W-2 forms for ALL EMPLOYEES who worked
for you in the previous year, run the Close For the Year program.
This program issues no requests, except for confirmation that it should
be run.  The program will not continue if the 940 and W-2 forms have
not been printed.
.he
.bp
.ig
CHAPTER 31
P R O D U C I N G   T H E   O P T I O N A L   R E P O R T S
.sp2
.pp5
The reports listed below can be printed at any time by choosing
the print program from the menu.  A 132 column printer
is assumed.  Be sure your printer is in the REMOTE mode if it
has a Local/Remote switch, and the paper aligned properly.  If
a printer malfunction should occur, press any key on the keyboard
to halt the printout.  To continue, type "RUN".  To end the print
program, type "STOP".
.he Operations: Chapter 31
.sp
.li
		   THE OPTIONAL REPORTS

		 1. ACCOUNT PRINT
		 2. EMPLOYEE MASTER REPORT
		 3. EMPLOYEE HISTORY REPORT
		 4. VACATION REPORT

	     (Figure 31.1: The Optional Reports)
.pp
The Account Print, Employee History Report,
and Vacation Report programs issue no requests.  Just align the paper
in your printer and select them from the menu.
.pp
The Employee Master List offers a full and a short version.
The full Employee Master Report prints one employee per page, the
short report one employee per line.  Answer the request for
which version you want by typing F for the full page version and
S for the single line version.
.pp
The Employee Master Report also offers a wide range of printing options.
These are presented in menu form (see Chapter 32 for a reproduction
of the menu).  Select the desired print option.  If you opt to print
a range of employees, the program asks you to enter the starting and
ending employee numbers.  The program checks your entries to make
sure the range is valid.
.he

